% Framework figure (Figure 1) - standalone TikZ
% Include in main paper with % Framework figure (Figure 1) - standalone TikZ
% Include in main paper with % Framework figure (Figure 1) - standalone TikZ
% Include in main paper with % Framework figure (Figure 1) - standalone TikZ
% Include in main paper with \input{fig_framework.tex}
\begin{figure}[t]
\centering
\resizebox{0.48\textwidth}{!}{%
\begin{tikzpicture}[
    node distance=0.6cm and 1.2cm,
    bluebox/.style={rectangle, draw=blue!60, fill=blue!10, rounded corners=3pt,
                minimum width=3.2cm, minimum height=0.8cm,
                font=\small, align=center, line width=0.8pt},
    greenbox/.style={rectangle, draw=green!50!black, fill=green!8, rounded corners=3pt,
                minimum width=3.2cm, minimum height=0.8cm,
                font=\small, align=center, line width=0.8pt},
    graybox/.style={rectangle, draw=gray!60, fill=gray!10, rounded corners=3pt,
                minimum width=3.2cm, minimum height=0.8cm,
                font=\small, align=center, line width=0.8pt},
    purplebox/.style={rectangle, draw=purple!60, fill=purple!10, rounded corners=3pt,
                minimum width=3.2cm, minimum height=0.8cm,
                font=\small, align=center, line width=0.8pt},
    bluephase/.style={rectangle, draw=blue!60, fill=blue!5, rounded corners=6pt,
                  minimum width=3.8cm, inner sep=10pt, line width=1.2pt},
    greenphase/.style={rectangle, draw=green!50!black, fill=green!5, rounded corners=6pt,
                  minimum width=3.8cm, inner sep=10pt, line width=1.2pt},
    bluearrow/.style={-{Stealth[length=5pt]}, line width=0.8pt, blue!60},
    greenarrow/.style={-{Stealth[length=5pt]}, line width=0.8pt, green!50!black},
    grayarrow/.style={-{Stealth[length=5pt]}, line width=0.8pt, gray!60},
    diamondnode/.style={diamond, draw=orange!80, fill=orange!10,
                    minimum width=1.2cm, minimum height=1.2cm,
                    font=\small, align=center, inner sep=1pt, line width=0.8pt},
    notelabel/.style={font=\footnotesize\itshape, text=gray!70!black},
]

% ===== TRAINING SECTION =====
\node[font=\large\bfseries] (title) {Training};

% Phase 1
\node[bluebox, below=0.5cm of title, xshift=-2.5cm] (data1) {Training Data\\{\footnotesize (Base Features)}};
\node[bluebox, below=0.5cm of data1] (train1) {Train Base Model\\{\footnotesize Bayesian optimization}};
\node[bluebox, below=0.5cm of train1] (save1) {Save Base Model\\{\footnotesize (model artifact)}};

% Phase 1 background
\begin{scope}[on background layer]
\node[bluephase, fit=(data1)(train1)(save1),
      label={[font=\small\bfseries, text=blue!70]above:Phase 1: Base Model}] (p1bg) {};
\end{scope}

% Phase 2
\node[greenbox, below=0.5cm of title, xshift=2.5cm] (load2) {Load Base Model\\{\footnotesize (artifact only)}};
\node[greenbox, below=0.5cm of load2] (train2) {Train with\\Extended Features\\{\footnotesize Adaptive tree pruning}};
\node[greenbox, below=0.5cm of train2] (save2) {Save Extended\\Model};

% Phase 2 background
\begin{scope}[on background layer]
\node[greenphase, fit=(load2)(train2)(save2),
      label={[font=\small\bfseries, text=green!50!black]above:Phase 2: Extended Model}] (p2bg) {};
\end{scope}

% Training arrows
\draw[bluearrow] (data1) -- (train1);
\draw[bluearrow] (train1) -- (save1);
\draw[greenarrow] (load2) -- (train2);
\draw[greenarrow] (train2) -- (save2);

% Transfer arrow
\draw[-{Stealth[length=5pt]}, line width=1pt, dashed, purple!70] (save1.east) -- ++(0.4,0) |- (load2.west)
    node[pos=0.25, above, notelabel] {model transfer};

% No data sharing annotation
\node[font=\scriptsize, text=red!70!black, align=center,
      below=0.1cm of p1bg.south east, xshift=0.8cm]
      {no raw data\\shared};

% ===== SEPARATOR =====
\draw[dashed, gray!50, line width=0.5pt]
    ([yshift=-0.4cm]p1bg.south west) -- ++(10.5,0);

% ===== INFERENCE SECTION =====
\node[font=\large\bfseries, below=1.2cm of save1, xshift=2.5cm] (inftitle) {Inference};

\node[graybox, below=0.4cm of inftitle] (newdata) {New Data};
\node[diamondnode, below=0.5cm of newdata] (decision) {\footnotesize Extended\\[-2pt]\footnotesize features?};

\node[bluebox, below left=0.6cm and 0.8cm of decision] (basemodel) {Base Model};
\node[greenbox, below right=0.6cm and 0.8cm of decision] (extmodel) {Extended Model};
\node[purplebox, below=0.6cm of decision, yshift=-0.8cm] (pred) {Prediction};

% Inference arrows
\draw[grayarrow] (newdata) -- (decision);
\draw[bluearrow] (decision.west) -| (basemodel)
    node[pos=0.25, above, font=\footnotesize] {No};
\draw[greenarrow] (decision.east) -| (extmodel)
    node[pos=0.25, above, font=\footnotesize] {Yes};
\draw[bluearrow] (basemodel) |- (pred);
\draw[greenarrow] (extmodel) |- (pred);

\end{tikzpicture}%
}
\caption{Population-aware two-phase framework. \textbf{Training:} Phase~1 builds a base model on commonly available features; Phase~2 loads only the model artifact (no raw data) and extends it with newly available features using adaptive tree pruning. \textbf{Inference:} data routes to the appropriate specialized model based on feature availability.}
\label{fig:framework}
\end{figure}

\begin{figure}[t]
\centering
\resizebox{0.48\textwidth}{!}{%
\begin{tikzpicture}[
    node distance=0.6cm and 1.2cm,
    bluebox/.style={rectangle, draw=blue!60, fill=blue!10, rounded corners=3pt,
                minimum width=3.2cm, minimum height=0.8cm,
                font=\small, align=center, line width=0.8pt},
    greenbox/.style={rectangle, draw=green!50!black, fill=green!8, rounded corners=3pt,
                minimum width=3.2cm, minimum height=0.8cm,
                font=\small, align=center, line width=0.8pt},
    graybox/.style={rectangle, draw=gray!60, fill=gray!10, rounded corners=3pt,
                minimum width=3.2cm, minimum height=0.8cm,
                font=\small, align=center, line width=0.8pt},
    purplebox/.style={rectangle, draw=purple!60, fill=purple!10, rounded corners=3pt,
                minimum width=3.2cm, minimum height=0.8cm,
                font=\small, align=center, line width=0.8pt},
    bluephase/.style={rectangle, draw=blue!60, fill=blue!5, rounded corners=6pt,
                  minimum width=3.8cm, inner sep=10pt, line width=1.2pt},
    greenphase/.style={rectangle, draw=green!50!black, fill=green!5, rounded corners=6pt,
                  minimum width=3.8cm, inner sep=10pt, line width=1.2pt},
    bluearrow/.style={-{Stealth[length=5pt]}, line width=0.8pt, blue!60},
    greenarrow/.style={-{Stealth[length=5pt]}, line width=0.8pt, green!50!black},
    grayarrow/.style={-{Stealth[length=5pt]}, line width=0.8pt, gray!60},
    diamondnode/.style={diamond, draw=orange!80, fill=orange!10,
                    minimum width=1.2cm, minimum height=1.2cm,
                    font=\small, align=center, inner sep=1pt, line width=0.8pt},
    notelabel/.style={font=\footnotesize\itshape, text=gray!70!black},
]

% ===== TRAINING SECTION =====
\node[font=\large\bfseries] (title) {Training};

% Phase 1
\node[bluebox, below=0.5cm of title, xshift=-2.5cm] (data1) {Training Data\\{\footnotesize (Base Features)}};
\node[bluebox, below=0.5cm of data1] (train1) {Train Base Model\\{\footnotesize Bayesian optimization}};
\node[bluebox, below=0.5cm of train1] (save1) {Save Base Model\\{\footnotesize (model artifact)}};

% Phase 1 background
\begin{scope}[on background layer]
\node[bluephase, fit=(data1)(train1)(save1),
      label={[font=\small\bfseries, text=blue!70]above:Phase 1: Base Model}] (p1bg) {};
\end{scope}

% Phase 2
\node[greenbox, below=0.5cm of title, xshift=2.5cm] (load2) {Load Base Model\\{\footnotesize (artifact only)}};
\node[greenbox, below=0.5cm of load2] (train2) {Train with\\Extended Features\\{\footnotesize Adaptive tree pruning}};
\node[greenbox, below=0.5cm of train2] (save2) {Save Extended\\Model};

% Phase 2 background
\begin{scope}[on background layer]
\node[greenphase, fit=(load2)(train2)(save2),
      label={[font=\small\bfseries, text=green!50!black]above:Phase 2: Extended Model}] (p2bg) {};
\end{scope}

% Training arrows
\draw[bluearrow] (data1) -- (train1);
\draw[bluearrow] (train1) -- (save1);
\draw[greenarrow] (load2) -- (train2);
\draw[greenarrow] (train2) -- (save2);

% Transfer arrow
\draw[-{Stealth[length=5pt]}, line width=1pt, dashed, purple!70] (save1.east) -- ++(0.4,0) |- (load2.west)
    node[pos=0.25, above, notelabel] {model transfer};

% No data sharing annotation
\node[font=\scriptsize, text=red!70!black, align=center,
      below=0.1cm of p1bg.south east, xshift=0.8cm]
      {no raw data\\shared};

% ===== SEPARATOR =====
\draw[dashed, gray!50, line width=0.5pt]
    ([yshift=-0.4cm]p1bg.south west) -- ++(10.5,0);

% ===== INFERENCE SECTION =====
\node[font=\large\bfseries, below=1.2cm of save1, xshift=2.5cm] (inftitle) {Inference};

\node[graybox, below=0.4cm of inftitle] (newdata) {New Data};
\node[diamondnode, below=0.5cm of newdata] (decision) {\footnotesize Extended\\[-2pt]\footnotesize features?};

\node[bluebox, below left=0.6cm and 0.8cm of decision] (basemodel) {Base Model};
\node[greenbox, below right=0.6cm and 0.8cm of decision] (extmodel) {Extended Model};
\node[purplebox, below=0.6cm of decision, yshift=-0.8cm] (pred) {Prediction};

% Inference arrows
\draw[grayarrow] (newdata) -- (decision);
\draw[bluearrow] (decision.west) -| (basemodel)
    node[pos=0.25, above, font=\footnotesize] {No};
\draw[greenarrow] (decision.east) -| (extmodel)
    node[pos=0.25, above, font=\footnotesize] {Yes};
\draw[bluearrow] (basemodel) |- (pred);
\draw[greenarrow] (extmodel) |- (pred);

\end{tikzpicture}%
}
\caption{Population-aware two-phase framework. \textbf{Training:} Phase~1 builds a base model on commonly available features; Phase~2 loads only the model artifact (no raw data) and extends it with newly available features using adaptive tree pruning. \textbf{Inference:} data routes to the appropriate specialized model based on feature availability.}
\label{fig:framework}
\end{figure}

\begin{figure}[t]
\centering
\resizebox{0.48\textwidth}{!}{%
\begin{tikzpicture}[
    node distance=0.6cm and 1.2cm,
    bluebox/.style={rectangle, draw=blue!60, fill=blue!10, rounded corners=3pt,
                minimum width=3.2cm, minimum height=0.8cm,
                font=\small, align=center, line width=0.8pt},
    greenbox/.style={rectangle, draw=green!50!black, fill=green!8, rounded corners=3pt,
                minimum width=3.2cm, minimum height=0.8cm,
                font=\small, align=center, line width=0.8pt},
    graybox/.style={rectangle, draw=gray!60, fill=gray!10, rounded corners=3pt,
                minimum width=3.2cm, minimum height=0.8cm,
                font=\small, align=center, line width=0.8pt},
    purplebox/.style={rectangle, draw=purple!60, fill=purple!10, rounded corners=3pt,
                minimum width=3.2cm, minimum height=0.8cm,
                font=\small, align=center, line width=0.8pt},
    bluephase/.style={rectangle, draw=blue!60, fill=blue!5, rounded corners=6pt,
                  minimum width=3.8cm, inner sep=10pt, line width=1.2pt},
    greenphase/.style={rectangle, draw=green!50!black, fill=green!5, rounded corners=6pt,
                  minimum width=3.8cm, inner sep=10pt, line width=1.2pt},
    bluearrow/.style={-{Stealth[length=5pt]}, line width=0.8pt, blue!60},
    greenarrow/.style={-{Stealth[length=5pt]}, line width=0.8pt, green!50!black},
    grayarrow/.style={-{Stealth[length=5pt]}, line width=0.8pt, gray!60},
    diamondnode/.style={diamond, draw=orange!80, fill=orange!10,
                    minimum width=1.2cm, minimum height=1.2cm,
                    font=\small, align=center, inner sep=1pt, line width=0.8pt},
    notelabel/.style={font=\footnotesize\itshape, text=gray!70!black},
]

% ===== TRAINING SECTION =====
\node[font=\large\bfseries] (title) {Training};

% Phase 1
\node[bluebox, below=0.5cm of title, xshift=-2.5cm] (data1) {Training Data\\{\footnotesize (Base Features)}};
\node[bluebox, below=0.5cm of data1] (train1) {Train Base Model\\{\footnotesize Bayesian optimization}};
\node[bluebox, below=0.5cm of train1] (save1) {Save Base Model\\{\footnotesize (model artifact)}};

% Phase 1 background
\begin{scope}[on background layer]
\node[bluephase, fit=(data1)(train1)(save1),
      label={[font=\small\bfseries, text=blue!70]above:Phase 1: Base Model}] (p1bg) {};
\end{scope}

% Phase 2
\node[greenbox, below=0.5cm of title, xshift=2.5cm] (load2) {Load Base Model\\{\footnotesize (artifact only)}};
\node[greenbox, below=0.5cm of load2] (train2) {Train with\\Extended Features\\{\footnotesize Adaptive tree pruning}};
\node[greenbox, below=0.5cm of train2] (save2) {Save Extended\\Model};

% Phase 2 background
\begin{scope}[on background layer]
\node[greenphase, fit=(load2)(train2)(save2),
      label={[font=\small\bfseries, text=green!50!black]above:Phase 2: Extended Model}] (p2bg) {};
\end{scope}

% Training arrows
\draw[bluearrow] (data1) -- (train1);
\draw[bluearrow] (train1) -- (save1);
\draw[greenarrow] (load2) -- (train2);
\draw[greenarrow] (train2) -- (save2);

% Transfer arrow
\draw[-{Stealth[length=5pt]}, line width=1pt, dashed, purple!70] (save1.east) -- ++(0.4,0) |- (load2.west)
    node[pos=0.25, above, notelabel] {model transfer};

% No data sharing annotation
\node[font=\scriptsize, text=red!70!black, align=center,
      below=0.1cm of p1bg.south east, xshift=0.8cm]
      {no raw data\\shared};

% ===== SEPARATOR =====
\draw[dashed, gray!50, line width=0.5pt]
    ([yshift=-0.4cm]p1bg.south west) -- ++(10.5,0);

% ===== INFERENCE SECTION =====
\node[font=\large\bfseries, below=1.2cm of save1, xshift=2.5cm] (inftitle) {Inference};

\node[graybox, below=0.4cm of inftitle] (newdata) {New Data};
\node[diamondnode, below=0.5cm of newdata] (decision) {\footnotesize Extended\\[-2pt]\footnotesize features?};

\node[bluebox, below left=0.6cm and 0.8cm of decision] (basemodel) {Base Model};
\node[greenbox, below right=0.6cm and 0.8cm of decision] (extmodel) {Extended Model};
\node[purplebox, below=0.6cm of decision, yshift=-0.8cm] (pred) {Prediction};

% Inference arrows
\draw[grayarrow] (newdata) -- (decision);
\draw[bluearrow] (decision.west) -| (basemodel)
    node[pos=0.25, above, font=\footnotesize] {No};
\draw[greenarrow] (decision.east) -| (extmodel)
    node[pos=0.25, above, font=\footnotesize] {Yes};
\draw[bluearrow] (basemodel) |- (pred);
\draw[greenarrow] (extmodel) |- (pred);

\end{tikzpicture}%
}
\caption{Population-aware two-phase framework. \textbf{Training:} Phase~1 builds a base model on commonly available features; Phase~2 loads only the model artifact (no raw data) and extends it with newly available features using adaptive tree pruning. \textbf{Inference:} data routes to the appropriate specialized model based on feature availability.}
\label{fig:framework}
\end{figure}

\begin{figure}[t]
\centering
\resizebox{0.48\textwidth}{!}{%
\begin{tikzpicture}[
    node distance=0.6cm and 1.2cm,
    bluebox/.style={rectangle, draw=blue!60, fill=blue!10, rounded corners=3pt,
                minimum width=3.2cm, minimum height=0.8cm,
                font=\small, align=center, line width=0.8pt},
    greenbox/.style={rectangle, draw=green!50!black, fill=green!8, rounded corners=3pt,
                minimum width=3.2cm, minimum height=0.8cm,
                font=\small, align=center, line width=0.8pt},
    graybox/.style={rectangle, draw=gray!60, fill=gray!10, rounded corners=3pt,
                minimum width=3.2cm, minimum height=0.8cm,
                font=\small, align=center, line width=0.8pt},
    purplebox/.style={rectangle, draw=purple!60, fill=purple!10, rounded corners=3pt,
                minimum width=3.2cm, minimum height=0.8cm,
                font=\small, align=center, line width=0.8pt},
    bluephase/.style={rectangle, draw=blue!60, fill=blue!5, rounded corners=6pt,
                  minimum width=3.8cm, inner sep=10pt, line width=1.2pt},
    greenphase/.style={rectangle, draw=green!50!black, fill=green!5, rounded corners=6pt,
                  minimum width=3.8cm, inner sep=10pt, line width=1.2pt},
    bluearrow/.style={-{Stealth[length=5pt]}, line width=0.8pt, blue!60},
    greenarrow/.style={-{Stealth[length=5pt]}, line width=0.8pt, green!50!black},
    grayarrow/.style={-{Stealth[length=5pt]}, line width=0.8pt, gray!60},
    diamondnode/.style={diamond, draw=orange!80, fill=orange!10,
                    minimum width=1.2cm, minimum height=1.2cm,
                    font=\small, align=center, inner sep=1pt, line width=0.8pt},
    notelabel/.style={font=\footnotesize\itshape, text=gray!70!black},
]

% ===== TRAINING SECTION =====
\node[font=\large\bfseries] (title) {Training};

% Phase 1
\node[bluebox, below=0.5cm of title, xshift=-2.5cm] (data1) {Training Data\\{\footnotesize (Base Features)}};
\node[bluebox, below=0.5cm of data1] (train1) {Train Base Model\\{\footnotesize Bayesian optimization}};
\node[bluebox, below=0.5cm of train1] (save1) {Save Base Model\\{\footnotesize (model artifact)}};

% Phase 1 background
\begin{scope}[on background layer]
\node[bluephase, fit=(data1)(train1)(save1),
      label={[font=\small\bfseries, text=blue!70]above:Phase 1: Base Model}] (p1bg) {};
\end{scope}

% Phase 2
\node[greenbox, below=0.5cm of title, xshift=2.5cm] (load2) {Load Base Model\\{\footnotesize (artifact only)}};
\node[greenbox, below=0.5cm of load2] (train2) {Train with\\Extended Features\\{\footnotesize Adaptive tree pruning}};
\node[greenbox, below=0.5cm of train2] (save2) {Save Extended\\Model};

% Phase 2 background
\begin{scope}[on background layer]
\node[greenphase, fit=(load2)(train2)(save2),
      label={[font=\small\bfseries, text=green!50!black]above:Phase 2: Extended Model}] (p2bg) {};
\end{scope}

% Training arrows
\draw[bluearrow] (data1) -- (train1);
\draw[bluearrow] (train1) -- (save1);
\draw[greenarrow] (load2) -- (train2);
\draw[greenarrow] (train2) -- (save2);

% Transfer arrow
\draw[-{Stealth[length=5pt]}, line width=1pt, dashed, purple!70] (save1.east) -- ++(0.4,0) |- (load2.west)
    node[pos=0.25, above, notelabel] {model transfer};

% No data sharing annotation
\node[font=\scriptsize, text=red!70!black, align=center,
      below=0.1cm of p1bg.south east, xshift=0.8cm]
      {no raw data\\shared};

% ===== SEPARATOR =====
\draw[dashed, gray!50, line width=0.5pt]
    ([yshift=-0.4cm]p1bg.south west) -- ++(10.5,0);

% ===== INFERENCE SECTION =====
\node[font=\large\bfseries, below=1.2cm of save1, xshift=2.5cm] (inftitle) {Inference};

\node[graybox, below=0.4cm of inftitle] (newdata) {New Data};
\node[diamondnode, below=0.5cm of newdata] (decision) {\footnotesize Extended\\[-2pt]\footnotesize features?};

\node[bluebox, below left=0.6cm and 0.8cm of decision] (basemodel) {Base Model};
\node[greenbox, below right=0.6cm and 0.8cm of decision] (extmodel) {Extended Model};
\node[purplebox, below=0.6cm of decision, yshift=-0.8cm] (pred) {Prediction};

% Inference arrows
\draw[grayarrow] (newdata) -- (decision);
\draw[bluearrow] (decision.west) -| (basemodel)
    node[pos=0.25, above, font=\footnotesize] {No};
\draw[greenarrow] (decision.east) -| (extmodel)
    node[pos=0.25, above, font=\footnotesize] {Yes};
\draw[bluearrow] (basemodel) |- (pred);
\draw[greenarrow] (extmodel) |- (pred);

\end{tikzpicture}%
}
\caption{Population-aware two-phase framework. \textbf{Training:} Phase~1 builds a base model on commonly available features; Phase~2 loads only the model artifact (no raw data) and extends it with newly available features using adaptive tree pruning. \textbf{Inference:} data routes to the appropriate specialized model based on feature availability.}
\label{fig:framework}
\end{figure}
